\chapter{总结与展望}


\section{总结}

本文基于非对称半桥反激式拓扑设计了一款副边反馈类型的AC/DC开关电源变换器芯片。随着信息时代的发展，在社会对便携式电子设备充电速度的需求下，中等输出功率的快充广泛应用。
非对称半桥反激式变换器由于其可实现全负载范围内零电压导通和低电压应力的优势被工业界重点关注。
为了提升非对称半桥反激式拓扑结构的优势，本文通过多种控制策略共同作用和改进现有控制方法实现高效控制，并设计关键技术进一步降低系统损耗，对克服现有瓶颈、提高变换器效率和功率密度有重要意义。

本文为了实现高输出功率提升应用领域，开展了一系列的研究。
首先针对非对称半桥反激式拓扑的工作原理进行了详细的分析。
接着为了实现变换器在全负载范围内的高效率控制，针对不同负载设计了多种控制策略；在空载工况下通过突发模式降低待机功耗；在轻载工况下将DCM模式改进为谐振谷值开关模式，实现高边功率管的零电压导通提高轻载效率；重载工况下采用CRM模式连续交替地导通功率管完成能量传递，维持稳定的高输出功率。

%接着设计了谷值锁定电路实现可靠的频率调控并且降低设计成本。

%最后为了进一步优化谐振谷值开关模式的效率和多模式下的高效率控制，分别提出了精确谷底导通和退磁时间动态校准技术方案。

本文基于SIMPLIS、Virtuoso电路设计平台和SMIC 0.18μm BCD工艺完成变换器芯片的电路设计、仿真和版图绘制。本文取得的主要工作内容和创新点如下：

1、基于SIMPLIS电路软件建立非对称半桥反激式电路模型，验证拓扑的工作原理和理论推导，完善电路的控制方式，为后续在Virtuoso软件中设计电路奠定基础。

2、针对宽输出电压电流范围下的高效率控制提出多模式控制、谷值锁定和带输出补偿的峰值电流控制技术。为了实现非对称半桥反激式变换器的高效控制，
在恒压控制模式中针对不同负载设计了不同的工作模式和脉宽调制方式，拟合最佳的负载频率曲线，负载突变时采用PWM+PFM调制方式同时控制峰值电流和开关频率变化实现高响应速度，稳定后维持不同模式下的峰值电流和开关频率的单一变化，防止系统振荡。
设计数模混合电路的谷值锁定方案，以更小的芯片面积功耗完成开关频率和谷值数量的锁定，避免在功率重叠范围内发生谐振波谷的跳谷不稳定现象。
通过带输出电压补偿的峰值电流控制电路匹配不同输出电压下负载电路和反馈电压的一一对应，控制反馈电压不随输出负载的变化而剧烈波动，影响多模式技术的错误控制。

3、为进一步降低零电压谐振谷值开关模式中功率管在谐振电压影响下的开关损耗，提出精确谷底导通技术。
%该技术电路工作于存在谐振电压的死区时间内，不同于传统反激式变换器的准谐振模式简单地使用过零比较器检测谷底导通下一开关周期，
该电路为了在谐振谷底精确导通功率管，实现最小开关损耗，不仅检测谷底位置还需考虑驱动电路的延时问题，故通过使用压控高通滤波器超前采样信号，旨在谐振谷底产生前产生逻辑控制信号，再采用锁相器模块将超前时间逐步逼近到等于驱动电路的延时时间，实现功率管栅极驱动信号恰好在谐振谷最低值导通功率管的功能。

4、提出退磁时间动态校准技术实现宽输出电压和全负载范围内的高效率控制。
低边功率管导通下的退磁时间内存在着导通损耗，退磁时间过长或过短都不利于传递效率的提高。过长的退磁时间增大导通损耗并可能影响下一周期副边二极管的反偏，产生可靠性问题；过短的导通时间无法将励磁能量传递完全，做不到零电流开关同样影响高效率控制。巧妙利用辅助绕组检测退磁完成信号降低电流检测成本，逐步实现退磁时间和能量传递时间的完全重叠，实现低边功率管的最佳导通时间控制。

\section{展望}

根据电路设计和仿真测试结果可知本文提到的变换器芯片性能良好，符合设计要求，存在多个关键模块优化系统效率。
但为了简化变换器的控制逻辑，副边只使用了续流二极管，存在一定的二极管整流损耗，未来考虑将加入同步整流技术更进一步地提高转换效率。
同时由于副边反馈结构因为仿真原因只能使用理想电路，因此最后的转换效率也较为理想。
因为作者本人的时间和能力问题，未进行流片测试，后续芯片流片测试后将测得更为精确的转换效率。





